\documentclass[11pt]{amsart}
\usepackage{amsmath,amssymb,amsthm}
\usepackage[margin=1in]{geometry}
\usepackage{booktabs}
\usepackage[hidelinks]{hyperref}

\newtheorem{theorem}{Theorem}[section]
\newtheorem{lemma}[theorem]{Lemma}
\newtheorem{proposition}[theorem]{Proposition}
\newtheorem{corollary}[theorem]{Corollary}
\newtheorem{conjecture}[theorem]{Conjecture}
\newtheorem*{theoremA}{Theorem A}
\newtheorem*{theoremB}{Theorem B}
\newtheorem*{theoremC}{Theorem C}
\newtheorem*{theoremD}{Theorem D}
\theoremstyle{definition}
\newtheorem{definition}[theorem]{Definition}
\theoremstyle{remark}
\newtheorem{remark}[theorem]{Remark}

\newcommand{\ord}{\operatorname{ord}}
\newcommand{\cO}{\mathcal{O}}
\newcommand{\Z}{\mathbb{Z}}
\newcommand{\Q}{\mathbb{Q}}
\newcommand{\F}{\mathbb{F}}
\DeclareMathOperator{\Nm}{N}

\title{Divisors of Mersenne numbers under doubling, and prime factors of $2^p-1$\\
beside near-Wagstaff numbers: sharp elementary bounds,\\ exact reductions, and barriers}
% Author line: replace with the actual author(s).
\author{[Author name]}
\date{November 2026}
\subjclass[2020]{Primary 11A51, 11N25; Secondary 11R18, 11B83}
\keywords{Mersenne numbers, divisor function, Wagstaff primes, cyclotomic polynomials, Erd\H{o}s problems}

\begin{document}

\begin{abstract}
Let $f(n)=\sum_{k\le n}\tau(2^k-1)$. Erd\H{o}s asked whether $f(2n)/f(n)$ tends to a limit
(Problem~\#893). Kova\v{c} and Luca proved $\limsup f(2n)/f(n)=\infty$. Whether the limit is $+\infty$ is open.
We prove that $f(2n)\ge 3f(n)$ for every $n$, with equality exactly for $n\in\{1,2,4\}$. If the Fermat number $F_a$ is
composite, then $f(2n)\ge\bigl(4-\tfrac{16}{3\cdot2^{a-2}+4}\bigr)f(n)$ on $[2^a,2^{a+1})$. Hence
$f(2n)\ge\tfrac{24}{7}f(n)$ for $32\le n<2^{33}$, and $\liminf f(2n)/f(n)\ge4$ if all large Fermat numbers are composite.
For every $C$, the set $\{n:f(2n)>Cf(n)\}$ meets every chain $n,2n,\dots,2^{J(C)-1}n$ and has lower density at least
$1/(2^{J(C)}-1)$. We reduce $f(2n)/f(n)\to\infty$ exactly to a statement about indices whose odd part has few
divisors. This gives a conditional theorem whose hypothesis is far weaker than that of Kova\v{c} and Luca.
We also show that one natural route to a negative answer needs primes $p$ for which
$\tau(2^p+1)$ is bounded and $\omega(2^p-1)=\exp\bigl((\log2+o(1))\log p/\log\log p\bigr)$.

This leads to the constant
$\beta_K=\limsup_{p:\,\tau(2^p+1)\le K}\log\omega(2^p-1)\,\log\log p/\log p$ and the conjecture $\beta_K=0$.
We do not settle it. We prove that a disproof would settle two open problems at once. We give barrier theorems, in
sketch form: axioms that capture the bounded-depth properties of the factorizations are consistent with $\beta_K=0$
and with $\beta_K=\infty$.
We give an ``honest ideal'' family showing that exact divisibility at the level of ideals is still not enough.
We prove the function-field analogue, which holds uniformly in $p$. All numerical claims were re-checked against independent
factor tables. This corrects a claim in the working notes: $\omega(2^{397}-1)=9$, not at most $8$.
\end{abstract}

\maketitle
\tableofcontents

%======================================================================
\section{Introduction}
%======================================================================

\subsection{Erd\H{o}s Problem \#893}
Write $\tau$ for the divisor function, $\omega$ for the number of distinct prime factors, and
\[
T(k)=\tau(2^k-1),\qquad f(n)=\sum_{k\le n}T(k).
\]
Erd\H{o}s~\cite{Er98} asked whether $f(2n)/f(n)$ tends to a limit; this is Problem~\#893 in~\cite{EP}.
He remarked that $f$ probably has no simple asymptotic formula because it grows too fast.
Kova\v{c} and Luca~\cite{KL} proved $\limsup_{n}f(2n)/f(n)=\infty$, so no finite limit exists. They also proved that
$f(2n)/f(n)\to\infty$ under either of two conjectures. One of them is
\begin{equation}\label{eq:KL}
\omega(\Phi_d(2))\le 10\log d\qquad\text{for all sufficiently large } d.
\end{equation}
Whether $f(2n)/f(n)\to\infty$ unconditionally, equivalently whether $\liminf f(2n)/f(n)=\infty$, is open.
To our knowledge no lower bound for the $\liminf$ beyond the trivial $f(2n)>2f(n)$ has been recorded.
Part~I proves the following.

\begin{theoremA}[Sharp constant; Theorems~\ref{thm:three}, \ref{thm:fermat}]
$f(2n)\ge3f(n)$ for all $n\ge1$, with equality exactly for $n\in\{1,2,4\}$. If $a\ge4$ and $F_a=2^{2^a}+1$ is
composite, then $f(2n)\ge\bigl(4-\frac{16}{3\cdot2^{a-2}+4}\bigr)f(n)$ for $2^a\le n<2^{a+1}$. In particular
$f(2n)\ge\frac{24}{7}f(n)$ for $32\le n<2^{33}$. If $F_a$ is composite for all large $a$, then
$\liminf f(2n)/f(n)\ge4$.
\end{theoremA}

\begin{theoremB}[Good $n$ are everywhere; Theorem~\ref{thm:density}]
Let $D(M)=\max_{m\le M}\tau(m)$ and $J(C)=\min\{J:D(2^J)-2>J\log_2C\}$. For every $C>1$ and every $n$, some
$0\le j<J(C)$ has $f(2^{j+1}n)>Cf(2^jn)$. Hence $\{n:f(2n)>Cf(n)\}$ has lower natural density at least
$1/(2^{J(C)}-1)$. Here $J(C)=O(\log\log C\cdot\log\log\log C)$.
\end{theoremB}

Theorem~B reproves the theorem of Kova\v{c} and Luca, since it gives $\limsup=\infty$.

\begin{theoremC}[Exact reduction; Theorem~\ref{thm:thin}, Corollary~\ref{cor:cond}]
Call $k$ \emph{$K$-thin} if the odd part of $k$ has fewer than $K$ divisors. Then $f(2n)/f(n)\to\infty$ if and only if,
for every fixed $K$, $\sum_{k\le n,\ k\ K\text{-thin}}T(k)=o\bigl(f(2n)-f(n)\bigr)$. Consequently
$f(2n)/f(n)\to\infty$ provided that, for every $K$, $\omega(\Phi_d(2))=o\bigl(\psi(d)/\log^2d\bigr)$ as $d\to\infty$
through $K$-thin $d$, where $\psi(x)=\max_{m\le x}\tau(m)=2^{(1+o(1))\log x/\log\log x}$.
\end{theoremC}

This hypothesis is much weaker than~\eqref{eq:KL}. It allows $\omega(\Phi_d(2))$ as large as $\psi(d)^{1-o(1)}$, and
it says nothing at all about indices that are not thin.

\begin{theoremD}[What a negative answer needs; Propositions~\ref{prop:negative}, \ref{prop:spike}]
If $f(2n)\le Cf(n)$ for infinitely many $n$, then for a fixed $K=K(C)$ there are $K$-thin
$d\le n$ with $\omega(\Phi_d(2))\ge(1-o(1))\psi(n)/(K\log_2^2n)$ for each such $n$. If moreover a single prime index $p\le n$ carries a fixed
fraction $1/\lambda$ of $f(n)$, then $\tau(2^p+1)\le C\lambda$ and
$\omega(2^p-1)\ge\exp\bigl((\log2+o(1))\log p/\log\log p\bigr)$ for infinitely many such $p$.
\end{theoremD}

\subsection{Mersenne numbers beside near-Wagstaff numbers}
Theorem~D isolates a concrete arithmetic scenario. Let $p$ be prime, $M_p=2^p-1$, and for $K\ge4$ let
\[
P_K=\{p:\ \tau(2^p+1)\le K\},\qquad
\beta_K=\limsup_{p\to\infty,\ p\in P_K}\frac{\log\omega(M_p)\cdot\log\log p}{\log p},
\]
with $\beta_K:=0$ if $P_K$ is finite. For $p>3$, $\tau(2^p+1)=4$ exactly when $(2^p+1)/3$ is a Wagstaff prime.
Theorem~D shows that $\beta_K\ge\log 2$ for some $K$ is necessary for the ``prime spike'' route to a negative answer
to Problem~\#893.

\begin{conjecture}\label{conj:beta}
$\beta_K=0$ for every $K\ge4$.
\end{conjecture}

Gillies' heuristic~\cite{Gi64} (see also Wagstaff~\cite{Wa83}) predicts that $\omega(M_p)$ is typically of
order $\log p$. This would make
Conjecture~\ref{conj:beta} true with room to spare. Unconditionally, nothing better than
$\omega(M_p)\le(\frac{\log2}{2}+o(1))\,p/\log p$ seems to be known pointwise. We do not resolve
Conjecture~\ref{conj:beta}. Part~II proves:
\begin{itemize}
\item a disproof would give infinitely many composite $M_p$ ($p$ prime) on which $(2^p+1)/3$ has at most $K/2-1$
prime factors (Theorem~\ref{thm:disproof});
\item barrier theorems: families of fake factorizations satisfy every bounded-depth property of the true ones, yet give
$\beta_K=0$ or $\beta_K=\infty$ as desired, and even saturate the trivial constant $\frac{\log2}2$
(Section~\ref{sec:barriers});
\item an ``honest ideal'' obstruction: elements $2^k-\zeta_p$ share every algebraic feature of $2-\zeta_p$ used by
class-group, Stickelberger, Galois and lattice methods, but have $\omega\ge\tau(k)$ (Theorem~\ref{thm:honest});
\item the function-field analogue: over $\F_\ell[T]$, $\tau(a^p+1)\le K$ implies
$\omega(a^p-1)\le\deg a\,(1+\log_2K)$, uniformly in $p$ (Theorem~\ref{thm:ff});
\item a divisor-lattice inequality and an incidence lemma that isolate the one estimate every method we tried needs
(Section~\ref{sec:tools}).
\end{itemize}

\subsection{What is new, and what is not}
Section~\ref{sec:novelty} compares each result with the literature. In short, \emph{neither problem is resolved here}.
The theorem $\limsup f(2n)/f(n)=\infty$ is due to Kova\v{c} and Luca. The trivial bounds, Zsigmondy-type facts and all factor
data are classical or tabulated. To our knowledge, the following are not in the literature: Theorems~A--D, the barrier
theorems, the honest-ideal obstruction, the function-field theorem and the divisor-lattice lemma.
All of them are elementary, or use only classical analytic input (Siegel--Walfisz, Bombieri--Vinogradov).
Several statements in the working notes behind this paper turned out to be false or imprecise.
They are corrected in Section~\ref{sec:errata}.

\begin{remark}[Provenance]
% Author: edit this remark as appropriate for the venue.
The results were first obtained in two working notes written with extensive help from AI research agents. For this
paper every proof given in full was re-derived and checked line by line. Every numerical claim was re-verified against
independent factor tables (Section~\ref{sec:data}). Results proved only in sketch form are labelled as such.
\end{remark}

%======================================================================
\section{Preliminaries}
%======================================================================
For odd $d$ let $o(d)=\ord_d2$. We use the following standard facts.

\begin{enumerate}
\item[(P1)] \emph{Cyclotomic factorization.} $2^k-1=\prod_{d\mid k}\Phi_d(2)$, and $2^m+1=\prod_{e\mid2m,\ e\nmid m}\Phi_e(2)$.
\item[(P2)] \emph{Bang--Zsigmondy}~\cite{Ba86,Zs92}. For $n\notin\{1,6\}$ there is a prime $q$ with $\ord_q2=n$
(a \emph{primitive} prime for $n$). Primitive primes for distinct $n$ are distinct. A primitive prime for $n$ divides
$2^k-1$ if and only if $n\mid k$.
\item[(P3)] \emph{Lifting the exponent.} If $q$ is an odd prime with $q\mid 2^m-1$, then
$v_q(2^{mt}-1)=v_q(2^m-1)+v_q(t)$.
\item[(P4)] \emph{Coprimality.} $\gcd(2^m-1,2^m+1)=1$, so $T(2m)=T(m)\,\tau(2^m+1)$.
\item[(P5)] \emph{Wigert}~\cite{Wi07}. $\psi(x):=\max_{m\le x}\tau(m)=2^{(1+o(1))\log x/\log\log x}$.
\end{enumerate}

\begin{lemma}[Divisor gain]\label{lem:gain}
Let $m\mid k$. Suppose $2^k-1$ has $s$ distinct prime factors that do not divide $2^m-1$. Then $T(k)\ge2^sT(m)$.
In particular, if $E$ is a set of divisors $e$ of $k$ with $e\nmid m$, then $T(k)\ge2^{|E\setminus\{1,6\}|}T(m)$.
\end{lemma}
\begin{proof}
$\tau(N)=\prod_{q\mid N}(v_q(N)+1)$. Every prime factor of $N'=2^m-1$ divides $N=2^k-1$ to at least the same power.
Each of the $s$ new primes contributes a factor $\ge2$. For the second claim apply (P2).
\end{proof}

\begin{lemma}\label{lem:Tupper}
$T(k)\le(k+1)^{\omega(2^k-1)}$ and $\omega(2^k-1)\le\tau(k)\max_{d\mid k}\omega(\Phi_d(2))$.
\end{lemma}
\begin{proof}
Each exponent satisfies $v_q(2^k-1)\le\log_3 2^k<k$. The second bound follows from (P1).
\end{proof}

\begin{lemma}\label{lem:plus}
(a) $\tau(2^m+1)\ge4$ unless $m$ is a power of $2$ or $m=3$; also $\tau(9)=3$.
(b) If $u$ is the odd part of $m$, then $\omega(2^m+1)\ge\tau(u)-1$.
\end{lemma}
\begin{proof}
(a) Suppose $m$ has an odd divisor $u>1$. Then $2^{m/u}+1$ is a proper divisor of $2^m+1$ greater than $1$.
If $\tau(2^m+1)=3$, then $2^m+1=q^2$, so $(q-1)(q+1)=2^m$, which forces $q=3$ and $m=3$.
(b) Write $m=2^cu$. By (P1) the indices $e\mid 2m$ with $e\nmid m$ are $e=2^{c+1}w$ with $w\mid u$. There are
$\tau(u)$ of them, and each except possibly $e=6$ has a primitive prime by (P2).
\end{proof}

%======================================================================
\part*{Part I. The doubling ratio}
\addcontentsline{toc}{part}{Part I. The doubling ratio}
%======================================================================

\section{An exact identity}

\begin{theorem}\label{thm:identity}
$f(n)=\sum_{d\ \mathrm{odd}}\lfloor n/o(d)\rfloor$. Hence
$f(2n)-2f(n)=\#\{d\ \mathrm{odd}:\{n/o(d)\}\ge\frac12\}$, where $\{\cdot\}$ is the fractional part.
\end{theorem}
\begin{proof}
An odd $d$ divides $2^k-1$ if and only if $o(d)\mid k$. So
$f(n)=\sum_{d}\#\{k\le n:o(d)\mid k\}=\sum_d\lfloor n/o(d)\rfloor$, where only finitely many $d$ contribute.
For real $x$, $\lfloor2x\rfloor-2\lfloor x\rfloor$ is $1$ if $\{x\}\ge\frac12$ and $0$ otherwise.
\end{proof}

The identity is elementary and probably folklore. It shows that $f(2n)>2f(n)$ always, since
$d=3$ gives $\{n/2\}=\frac12$ for odd $n$, and $d=2^{n+1}-1$ gives $o(d)=n+1>n$.
It also recasts the problem: do the moduli $d$ whose order puts $n$ in the upper half of a period outnumber
$\sum_d\lfloor n/o(d)\rfloor$ by an unbounded factor?

\section{The sharp constant 3}

Write $F_i=2^{2^i}+1$. By (P4), $T(2^{a+1})=\tau(F_a)\,T(2^a)$ and $T(2^a)=\prod_{i<a}\tau(F_i)$.

\begin{lemma}\label{lem:Cr}
Let $r=2^bu$ with $u>1$ odd, and $C(r)=\sum_{i=0}^{b}T(2^iu)$. Then $C(r)\le T(2r)/3$.
\end{lemma}
\begin{proof}
By (P4), $T(2^{i+1}u)=T(2^iu)\,\tau(2^{2^iu}+1)$. Since $u>1$, the exponent $2^iu$ is never a power of $2$.
By Lemma~\ref{lem:plus}(a), each factor $\tau(2^{2^iu}+1)$ is $\ge4$, except the factor $\tau(9)=3$ when $i=0$, $u=3$.
If all factors are $\ge4$, then $T(2^iu)\le4^{-(b+1-i)}T(2r)$, and $C(r)\le T(2r)\sum_{j\ge1}4^{-j}=T(2r)/3$.
If $u=3$, the same bound holds for $i\ge1$, and $T(3)=T(6)/3\le T(2r)/(3\cdot4^b)$. Then
$C(r)\le T(2r)\bigl(\sum_{j=1}^{b}4^{-j}+\frac{1}{3\cdot4^b}\bigr)=T(2r)/3$.
\end{proof}

\begin{lemma}\label{lem:3pow}
For $c\ge1$, $T(3\cdot2^c)\ge\frac32\,2^c\,T(2^c)$.
\end{lemma}
\begin{proof}
Let $N=2^{3\cdot2^c}-1$ and $N'=2^{2^c}-1$. The indices $3\cdot2^i$, $0\le i\le c$, $i\ne1$, divide $3\cdot2^c$ but not
$2^c$. Their $c$ primitive primes (P2) are distinct, do not divide $N'$, and differ from $3$, because $\ord_32=2$.
By (P3), $v_3(N')=v_3(4^{2^{c-1}}-1)=1$ and $v_3(N)=1+v_3(3\cdot2^{c-1})=2$. So the prime $3$ contributes
$3$ to $\tau(N)$ against $2$ to $\tau(N')$, and $\tau(N)\ge\tau(N')\cdot\frac32\cdot2^c$.
\end{proof}

\begin{lemma}\label{lem:fermatsum}
$T(2^{i+1})\ge2T(2^i)$. Hence $\sum_{i<a}T(2^i)<2T(2^{a-1})$ for $a\ge1$.
\end{lemma}
\begin{proof}
$\tau(F_i)\ge2$. Then $\sum_{i<a}T(2^i)\le T(2^{a-1})\sum_{j=0}^{a-1}2^{-j}<2T(2^{a-1})$.
\end{proof}

\paragraph{The decomposition.}
Let $R=\Z\cap(n/2,n]$. Every $m\le n$ is uniquely $m=r/2^j$ with $r\in R$ and $0\le j\le v_2(r)$. Hence
\begin{equation}\label{eq:decomp}
f(n)=\sum_{r\in R}C(r),\qquad C(r)=\sum_{j=0}^{v_2(r)}T(r/2^j),
\end{equation}
which agrees with Lemma~\ref{lem:Cr} for $r=2^bu$. The numbers $2r$, $r\in R$, are distinct elements of $(n,2n]$.
So $f(2n)\ge f(n)+\sum_{r\in R}T(2r)$, and for every $\lambda$,
\begin{equation}\label{eq:master}
f(2n)-(\lambda+1)f(n)\ \ge\ \sum_{r\in R}\bigl(T(2r)-\lambda\,C(r)\bigr).
\end{equation}

\begin{theorem}\label{thm:three}
$f(2n)\ge3f(n)$ for every $n\ge1$, with equality exactly for $n\in\{1,2,4\}$.
\end{theorem}
\begin{proof}
For $n\le15$ this is a direct computation from
$T(1),\dots,T(30)=1,2,2,4,2,6,2,8,4,8,4,24,2,8,8,16,2,32,2,48,12,16,4,96,8,8,8,64,8,96$.
The ratios $f(2n)/f(n)$ for $n=1,\dots,15$ are
$3,\,3,\,3.4,\,3,\,3.545,\,3.941,\,4.053,\,3.741,\,4.355,\,4.744,\,4.953,\,4.672,\,4.768,\,5.208,\,5.941$.

Let $n\ge16$ and $2^a\le n<2^{a+1}$, so $a\ge4$. Apply~\eqref{eq:master} with $\lambda=2$. For $r\in R$, $r\ne2^a$,
Lemma~\ref{lem:Cr} gives $T(2r)-2C(r)\ge T(2r)/3\ge0$. For $r=2^a$, Lemma~\ref{lem:fermatsum} gives
\[
T(2^{a+1})-2\sum_{i\le a}T(2^i)>(\tau(F_a)-2)T(2^a)-4T(2^{a-1})\ge-4T(2^{a-1}).
\]
To compensate, let $r_0=3\cdot2^{a-2}$ if $n<3\cdot2^{a-1}$, and $r_0=3\cdot2^{a-1}$ otherwise. In both cases
$r_0\in R$ and $r_0\ne2^a$. Also $2^{3\cdot2^{a-1}}-1$ divides $2^{2r_0}-1$, so
$T(2r_0)\ge T(3\cdot2^{a-1})$. By Lemma~\ref{lem:3pow},
$\tfrac13T(2r_0)\ge\tfrac13T(3\cdot 2^{a-1})\ge2^{a-2}T(2^{a-1})\ge4T(2^{a-1})$.
The sum in~\eqref{eq:master} is therefore strictly positive.
\end{proof}

\begin{remark}[Sharpness for soft methods]\label{rem:sharp3}
The constant $3$ cannot be improved by any argument that uses only the soft properties listed in
Section~\ref{sec:barrier-doubling}. Consider a model in which $\Phi_{2^a}(2)=F_{a-1}$ is replaced by a product of
many primes, so that $T^*(2^a)$ dominates $f^*(2^a)$, and $F_a$ is ``prime''. Then
$T^*(2^{a+1})=2T^*(2^a)$, so $f^*(2^{a+1})/f^*(2^a)\to3$.
\end{remark}

\begin{theorem}\label{thm:fermat}
If $a\ge4$ and $F_a$ is composite, then $f(2n)\ge\bigl(4-\frac{16}{3\cdot2^{a-2}+4}\bigr)f(n)$ for all
$n\in[2^a,2^{a+1})$.
\end{theorem}
\begin{proof}
Apply~\eqref{eq:master} with $\lambda=3$. The terms with $r\ne2^a$ are $\ge0$ by Lemma~\ref{lem:Cr}.
Since $F_a\equiv2\pmod3$ is not a square, compositeness gives $\tau(F_a)\ge4$. With Lemma~\ref{lem:fermatsum}, the
$r=2^a$ term is then $\ge T(2^a)-6T(2^{a-1})\ge-4T(2^{a-1})$. Put $t=T(2^{a-1})/f(n)$; this gives
$f(2n)/f(n)\ge4-4t$. On the other hand $3\cdot2^{a-1}\le2n$, so by Lemma~\ref{lem:3pow}
$f(2n)\ge T(3\cdot2^{a-1})\ge3\cdot2^{a-2}\,T(2^{a-1})$, that is, $f(2n)/f(n)\ge3\cdot2^{a-2}t$. The larger of the two
bounds is smallest where they meet, at $t=4/(3\cdot2^{a-2}+4)$.
\end{proof}

\begin{corollary}\label{cor:fermat}
(a) $F_5,\dots,F_{32}$ are composite, so $f(2n)\ge\frac{24}{7}f(n)$ for $32\le n<2^{33}$.
(b) If $F_a$ is composite for all sufficiently large $a$, then $\liminf_{n}f(2n)/f(n)\ge4$.
\end{corollary}
\begin{proof}
(a) Take $a=5$ in Theorem~\ref{thm:fermat}; the bound increases with $a$. (b) Every large $n$ lies in some
$[2^a,2^{a+1})$, and the bound tends to $4$ as $a\to\infty$.
\end{proof}

The value $4$ in (b) is the value predicted by the naive ``random integer'' model, which Kova\v{c} and Luca show to be misleading.
Exact data give $\min_{32\le n\le606}f(2n)/f(n)=14.26$ (Section~\ref{sec:data}).

\section{Growth along chains and the density of good \texorpdfstring{$n$}{n}}

\begin{proposition}\label{prop:mult}
For all $M,n\ge1$, $f(Mn)\ge2^{D(M)-2}f(n)$, where $D(M)=\max_{m\le M}\tau(m)$.
\end{proposition}
\begin{proof}
Choose $h\le M$ with $\tau(h)=D(M)$, and fix $m\le n$. The indices $me$ with $e\mid h$, $e>1$, divide $hm$ but not $m$.
There are $\tau(h)-1$ of them, and at most one equals $6$. By Lemma~\ref{lem:gain}, $T(hm)\ge2^{\tau(h)-2}T(m)$.
The numbers $hm$, $m\le n$, are distinct and at most $Mn$. Summing over $m$ gives the claim.
\end{proof}

\begin{theorem}\label{thm:density}
Let $J=J(C)=\min\{J:D(2^J)-2>J\log_2C\}$. For every $n$ some $0\le j<J$ satisfies $f(2^{j+1}n)>Cf(2^jn)$.
Consequently $\{n:f(2n)>Cf(n)\}$ has lower natural density at least $1/(2^J-1)$. This bound is optimal for the chain
property alone.
\end{theorem}
\begin{proof}
By Proposition~\ref{prop:mult} with $M=2^J$, $\prod_{j<J}f(2^{j+1}n)/f(2^jn)=f(2^Jn)/f(n)\ge2^{D(2^J)-2}>C^J$.
So some factor exceeds $C$.

For the density, group $n\le x$ by odd part $m$. The chain $m,2m,4m,\dots$ has
$\lfloor\log_2(x/m)\rfloor+1$ terms up to $x$. Every $J$ consecutive terms contain a good one, so at least
$\lfloor(\lfloor\log_2(x/m)\rfloor+1)/J\rfloor=\#\{i\ge0:m2^{iJ+J-1}\le x\}$ terms are good. Summing over odd $m$ gives
$\sum_{i\ge0}\frac{x}{2^{iJ+J}}+O(J\log x)=\frac{x}{2^J-1}+O(J\log x)$.
The set $\{n:v_2(n)\equiv J-1\pmod J\}$ meets every chain of $J$ consecutive terms and has density exactly
$1/(2^J-1)$.
\end{proof}

By direct computation, $D(2^J)$ for $J=7,8,12,13,20,21$ equals $16,20,48,64,240,288$. Hence $J(4)=8$, $J(16)=13$ and
$J(10^4)=21$. By (P5), $\log_2D(2^J)=(1+o(1))J\log2/\log J$, so $J(C)=O(\log\log C\cdot\log\log\log C)$.

\begin{corollary}[Kova\v{c}--Luca~\cite{KL}]\label{cor:KL}
$\limsup_nf(2n)/f(n)=\infty$. More strongly, $f(M(n)n)/f(n)\to\infty$ for every integer sequence $M(n)\to\infty$.
\end{corollary}

The argument of~\cite{KL} also uses highly composite numbers, through $f'(n)=\sum_{k\le n}2^{\tau(k)}$.
Theorem~\ref{thm:density} adds the chain and density statements.

\section{Exact reduction to thin indices}

\begin{definition}
Fix $K\ge2$. An integer $k$ is \emph{$K$-thin} if its odd part has fewer than $K$ divisors. Every divisor of a
$K$-thin number is $K$-thin.
\end{definition}

\begin{theorem}\label{thm:thin}
For all $n$ and all $K\ge2$,
\[
\sum_{\substack{k\le n\\ k\text{ not }K\text{-thin}}}T(k)\ \le\ \frac{f(2n)-f(n)}{2^{K-1}-1}.
\]
Hence $f(2n)/f(n)\to\infty$ if and only if, for every fixed $K$, the $T$-mass of the $K$-thin $k\le n$ is
$o\bigl(f(2n)-f(n)\bigr)$.
\end{theorem}
\begin{proof}
For $k\le n$ there is a unique $b\ge1$ with $N=2^bk\in(n,2n]$. The map $k\mapsto(N,b)$ is a bijection onto
$\{(N,b):N\in(n,2n],\ 1\le b\le v_2(N)\}$. If $k$ is not $K$-thin, then every $2^ik$ has the same odd part $u$ with
$\tau(u)\ge K$. By (P4) and Lemma~\ref{lem:plus}(b), each doubling multiplies $T$ by at least $2^{K-1}$, so
$T(N)\ge2^{b(K-1)}T(k)$. Therefore
\[
\sum_{k\text{ not thin}}T(k)\le\sum_{N\in(n,2n]}T(N)\sum_{b\ge1}2^{-b(K-1)}=\frac{f(2n)-f(n)}{2^{K-1}-1}.
\]
For the equivalence: if the thin mass is $o(f(2n)-f(n))$ for every $K$, then
$\limsup f(n)/(f(2n)-f(n))\le1/(2^{K-1}-1)$ for every $K$, so the ratio tends to infinity. The converse is trivial,
since the thin mass is at most $f(n)$.
\end{proof}

\begin{corollary}[Conditional divergence]\label{cor:cond}
Suppose that for every $K\ge2$,
\begin{equation}\label{eq:thinhyp}
\omega(\Phi_d(2))=o\!\left(\frac{\psi(d)}{\log^2d}\right)\qquad(d\to\infty,\ d\ K\text{-thin}).
\end{equation}
Then $f(2n)/f(n)\to\infty$. In particular this holds under~\eqref{eq:KL}, or under
$\omega(\Phi_d(2))\le\psi(d)^{1-\delta}$ for some $\delta>0$ and all large $d$.
\end{corollary}
\begin{proof}
Fix $K$ and a $K$-thin $k\le n$. Then $\tau(k)<K(\log_2n+1)$, and every divisor $d$ of $k$ is $K$-thin. We claim that
$\max_{d\mid k}\omega(\Phi_d(2))=o(\psi(n)/\log^2n)$, uniformly in $k\le n$.
Bounded $d$ contribute $O_K(1)$. If $d\ge\sqrt n$, then $\log d\ge\frac12\log n$ and $\psi(d)\le\psi(n)$, so
\eqref{eq:thinhyp} gives the claim, uniformly because $d\ge\sqrt n\to\infty$. If $d<\sqrt n$, then
\eqref{eq:thinhyp} gives $\omega(\Phi_d(2))\le\psi(d)+O_K(1)\le\psi(\sqrt n)+O_K(1)=\psi(n)^{1/2+o(1)}$. This is
$o(\psi(n)/\log^2n)$ because $\psi(n)$ grows faster than any power of $\log n$.
By Lemma~\ref{lem:Tupper},
\[
\log_2T(k)\le\log_2(n+1)\cdot K(\log_2n+1)\cdot o\bigl(\psi(n)/\log^2n\bigr)=o(\psi(n)).
\]
So the thin mass is at most $n\,2^{o(\psi(n))}=2^{o(\psi(n))}$. On the other hand, take $m^*\le n$ with
$\tau(m^*)=\psi(n)$ and let $N$ be the multiple $2^jm^*$ that lies in $(n,2n]$. Since $2^{m^*}-1\mid2^N-1$, and by
(P2), $f(2n)-f(n)\ge T(N)\ge T(m^*)\ge2^{\psi(n)-2}$. Theorem~\ref{thm:thin} now applies. The two particular cases imply
\eqref{eq:thinhyp} because $\log d$ and $\psi(d)^{1-\delta}$ are both $o(\psi(d)/\log^2d)$.
\end{proof}

\begin{proposition}[What a negative answer needs]\label{prop:negative}
Suppose $f(2n)\le Cf(n)$ for infinitely many $n$, and fix $K$ with $2^{K-1}-1>2C$. Then for each such $n$ there are a
$K$-thin $k\le n$ with $\log_2T(k)\ge\psi(n)-3-\log_2n$, and a $K$-thin $d\mid k$ with
\[
\omega(\Phi_d(2))\ \ge\ \frac{\psi(n)-3-\log_2n}{K(\log_2n+1)\log_2(n+1)}.
\]
\end{proposition}
\begin{proof}
By Theorem~\ref{thm:thin}, the non-thin mass is at most $(f(2n)-f(n))/(2^{K-1}-1)<f(n)/2$. So the thin mass is at
least $f(n)/2\ge2^{\psi(n)-3}$, using $f(n)\ge T(m^*)\ge2^{\psi(n)-2}$. There are at most $n$ thin $k\le n$,
which gives the first claim. The second follows from Lemma~\ref{lem:Tupper} and $\tau(k)<K(\log_2n+1)$.
\end{proof}

\begin{proposition}[Prime spikes force near-Wagstaff Mersenne numbers]\label{prop:spike}
Let $C,\lambda\ge1$. Suppose there are infinitely many $n$ with $f(2n)\le Cf(n)$ for which some prime $p\le n$
satisfies $T(p)\ge f(n)/\lambda$. Then $P_{C\lambda}$ is infinite and $\beta_{C\lambda}\ge\log2$.
\end{proposition}
\begin{proof}
Since $2p\le2n$, (P4) gives $\tau(2^p+1)=T(2p)/T(p)\le f(2n)/T(p)\le C\lambda$, so $p\in P_{C\lambda}$. As
$n\to\infty$ we have $T(p)\ge f(n)/\lambda\to\infty$, so infinitely many distinct $p$ occur. Next,
$\log_2T(p)\ge\log_2f(n)-\log_2\lambda\ge\psi(n)-2-\log_2\lambda\ge\psi(p)-2-\log_2\lambda$. By
Lemma~\ref{lem:Tupper}, $\omega(M_p)\ge(\psi(p)-2-\log_2\lambda)/\log_2(p+1)$. Hence
$\log\omega(M_p)\ge(\log2+o(1))\log p/\log\log p$ by (P5).
\end{proof}

So Conjecture~\ref{conj:beta} would rule out the simplest mechanism for a bounded ratio: a prime index carrying a
fixed share of $f(n)$. The soft-property models of Theorem~\ref{thm:model-doubling} realize exactly this mechanism.

%======================================================================
\part*{Part II. Prime factors of \texorpdfstring{$2^p-1$}{2\textasciicircum p-1} when \texorpdfstring{$2^p+1$}{2\textasciicircum p+1} has few divisors}
\addcontentsline{toc}{part}{Part II. Prime factors of \texorpdfstring{$2^p-1$}{2\textasciicircum p-1} when \texorpdfstring{$2^p+1$}{2\textasciicircum p+1} has few divisors}
%======================================================================

\section{Elementary structure}
Throughout, $p\ge5$ is prime, $a_p=\Phi_{2p}(2)=(2^p+1)/3$, $\zeta=\zeta_p$ and $\cO=\Z[\zeta]$.

\begin{lemma}\label{lem:3part}
$3\parallel2^p+1$, so $\tau(2^p+1)=2\tau(a_p)$.
\end{lemma}
\begin{proof}
$\ord_92=6$ and $p$ is odd and prime to $3$, so $2^p+1\equiv3$ or $6\pmod9$.
\end{proof}

\begin{lemma}\label{lem:mod8}
(a) $a_p\equiv3\pmod8$. (b) Every prime $Q\mid a_p$ has $\ord_Q2=2p$, so $Q\equiv1\pmod{2p}$.
(c) $Q\equiv3\pmod8$ if $(Q-1)/2p$ is odd, and $Q\equiv1\pmod8$ otherwise.
(d) The number of prime factors $\equiv3\pmod8$ of $a_p$, counted with multiplicity, is odd.
\end{lemma}
\begin{proof}
$2^p\equiv8\pmod{24}$ gives (a). For (b): $2^p\equiv-1\pmod Q$ and $Q\ne3$ by Lemma~\ref{lem:3part}.
For (c), Euler's criterion gives $(\frac2Q)\equiv(2^p)^{(Q-1)/2p}\equiv(-1)^{(Q-1)/2p}$. Combine this with the parity
of $(Q-1)/2=p\cdot\frac{Q-1}{2p}$, which fixes $Q\bmod4$. (d) follows from (a), since $3^k\equiv3\pmod8$ if and only
if $k$ is odd.
\end{proof}

\begin{corollary}\label{cor:4div}
$a_p$ is never a square. Hence $4\mid\tau(2^p+1)$ and $P_K=P_{4\lfloor K/4\rfloor}$. In particular $P_6=P_7=P_4$.
\end{corollary}

\begin{lemma}[Disjoint fibres]\label{lem:fibres}
Prime factors $q$ of $M_p$ have $\ord_q2=p$, and prime factors $Q$ of $a_p$ have $\ord_Q2=2p$. So no prime divides
two members of $\{M_p\}\cup\{a_p\}$, and no covering-congruence argument can force $a_p$ to be composite.
\end{lemma}

\begin{lemma}[Trivial bound]\label{lem:trivial}
$\omega(M_p)\le(\frac{\log2}2+o(1))\,p/\log p$.
\end{lemma}
\begin{proof}
Let $q_1<\dots<q_k$ be the prime factors. They are $\equiv1\pmod{2p}$, so $q_j>2pj$ and $(2p)^kk!<2^p$. We may
assume $k\ge p/\log^2p$. Then $\log k\ge\log p-2\log\log p$ and $k!\ge(k/e)^k$ give
$k(2\log p-2\log\log p+\log 2-1)<p\log2$.
\end{proof}

Using that $q$ lies in two classes modulo $8p$ and the Montgomery--Vaughan form of the Brun--Titchmarsh
inequality~\cite{MV73}, one gets the second-order form $\omega(M_p)\le\frac{\log2}2\,p/(\log p-1.0298+o(1))$.
The constant $\frac{\log2}{2}$ is unchanged.

\begin{remark}[Scale of the problem]
The quantity in $\beta_K$ lies in $[0,(1+o(1))\log\log p]$. So even $\beta_K<\infty$ requires
$\omega(M_p)\le p^{O(1/\log\log p)}$ on $P_K$, far beyond Lemma~\ref{lem:trivial}. At most
$p\log2/\log X$ prime factors exceed $X$. So the range that must be controlled is
$q\in[p^{1+o(1)},\,2^{p^{1-o(1)}}]$, not a range with a fixed $\varepsilon$.
\end{remark}

\section{What a disproof, or finiteness, would mean}

\begin{theorem}\label{thm:disproof}
Suppose $\beta_K>0$ and $0<c<\beta_K$. Then there are infinitely many primes $p$ such that
$\Omega(a_p)\le K/2-1$, $\omega(a_p)\le\log_2(K/2)$, and $\omega(M_p)\ge\exp(c\log p/\log\log p)$. In particular
these $M_p$ are composite.
\end{theorem}
\begin{proof}
$\beta_K>0$ forces $P_K$ to be infinite. Apply $\tau(n)\ge\Omega(n)+1$ and $\tau(n)\ge2^{\omega(n)}$ to
$\tau(a_p)\le K/2$.
\end{proof}

The infinitude of composite $M_p$ with $p$ prime is open (Guy~\cite{Guy}, A3). A disproof would settle it together
with an almost-prime statement for $(2^p+1)/3$ on the same set of $p$.

\begin{proposition}\label{prop:finite}
$P_K$ is finite for every $K$ if and only if $\Omega(a_p)\to\infty$. In particular, finiteness of $P_4$ means that
there are only finitely many Wagstaff primes.
\end{proposition}
This is believed false. The Lenstra--Pomerance--Wagstaff-type heuristic (see~\cite{BSW89}) predicts
$\asymp\log x$ Wagstaff exponents $p\le x$.

\begin{proposition}[No fixed-family certificate]\label{prop:nocert}
Let $f_1,\dots,f_r\in\Z[x]$ be nonconstant. There is a set of primes $p$ of positive relative density on which no
$f_j(p)$ is a divisor of $a_p$ other than $\pm1$, and likewise for $M_p$.
\end{proposition}
\begin{proof}
If $f_j=c\,x^k$, then $f_j(p)=cp^k$ is never such a divisor, because every prime factor of $a_pM_p$ is at least
$2p+1$. Otherwise write $f_j=x^kg_j$ with $g_j(0)\ne0$ and $g_j$ nonconstant. By Schur's theorem infinitely many primes
$\ell$ divide some value $g_j(r)$. Choose such odd primes $\ell_j$, distinct and not dividing $g_j(0)$, with residues
$r_j$. Then $r_j\not\equiv0\pmod{\ell_j}$. By the Chinese remainder theorem and Dirichlet's theorem, a positive proportion of primes satisfy
$p\equiv r_j\pmod{\ell_j}$ for all $j$. For these $p$ we have $\ell_j\mid f_j(p)$, while $\ell_j<2p+1$ for large $p$. So
$f_j(p)$ divides neither $a_p$ nor $M_p$, unless $f_j(p)=0$, which happens for finitely many $p$.
\end{proof}

The working note's proof of Proposition~\ref{prop:nocert} omitted the case $f_j=cx^k$, in which no
nonzero residue works; the repair is the first sentence above.

%======================================================================
\section{Barriers for soft and bounded-depth arguments}\label{sec:barriers}
%======================================================================
The results in this section are \emph{meta-theorems}. Each lists properties of the true factorizations and builds
fake factorizations with all of them but with opposite behaviour. Any argument that uses only the listed properties
therefore proves nothing about the question. The constructions are explicit, but they rely on the
Siegel--Walfisz or Bombieri--Vinogradov theorems and are ineffective. Full verifications are in the working notes.
Here we give statements and the key steps; these results are \emph{proved in sketch form only}.

\subsection{Soft properties of \texorpdfstring{$2^k-1$}{2\textasciicircum k-1}}\label{sec:barrier-doubling}
The \emph{soft properties} are the following:
\begin{itemize}
\item the cyclotomic factorization and $\gcd(2^a-1,2^b-1)=2^{\gcd(a,b)}-1$;
\item a primitive prime $\equiv1\pmod d$ for every index $d\ne1,6$;
\item $T(2k)=T(k)\tau(2^k+1)$ and lifting-the-exponent multiplicities;
\item the size $\log\Phi_d(2)\approx\varphi(d)\log 2$ and Aurifeuillian splitting;
\item the rule $q\equiv\pm1\pmod8$ if and only if $2\mid(q-1)/\ord_q2$, and power-residue rules of bounded depth;
\item Stewart's lower bound for the largest prime factor~\cite{St13};
\item Brun--Titchmarsh and averaged sieve bounds.
\end{itemize}

\begin{theorem}[Models for the doubling problem; sketch]\label{thm:model-doubling}
There are ``fake'' factorizations, in which each $\Phi_d(2)$ is replaced by a product of genuine primes
$\equiv1\pmod d$ satisfying all soft properties, of the following kind. They have one prime per generic index (two at
Aurifeuillian indices). Along a sparse sequence of primes $\ell_j$, index $\ell_j$ carries
$W_j\ge D(3\ell_j)+O(\log\ell_j)$ primes and index $2\ell_j$ carries one. For these models
$f^*(2\ell_j)/f^*(\ell_j)\to5$, so $\liminf f^*(2n)/f^*(n)<\infty$. Other models, with every
$\omega(\Phi_d^*(2))\le d^{o(1)}$, have $f^*(2n)/f^*(n)\to\infty$.
\end{theorem}
\begin{proof}[Key steps]
In $(\ell_j,2\ell_j]$ the only multiples of $\ell_j$ are $\ell_j$ and $2\ell_j$, so
$T^*(\ell_j)=2^{W_j}$ and $T^*(2\ell_j)=4T^*(\ell_j)$. Every other $k\le2\ell_j$ has
$\log_2T^*(k)\le D(2\ell_j)+\sum_{i<j}W_i+O(\log \ell_j)$. When the $\ell_j$ grow fast, this is negligible against $W_j$.
Then $f^*(\ell_j)\sim T^*(\ell_j)$ and $f^*(2\ell_j)\sim5T^*(\ell_j)$. Size is compatible because
$W_j\ll\ell_j/\log\ell_j$ primes $\equiv1\pmod{\ell_j}$ fit inside $2^{\ell_j}$. Primes in the required residue
classes, of modulus $O(d)$, exist by Siegel--Walfisz.
\end{proof}

The model at index $2\ell_j$ has $\Phi^*_{2\ell_j}(2)$ ``prime'', that is, $\tau(2^{\ell_j}+1)=4$, and
$\omega(M^*_{\ell_j})$ is huge. This is exactly the near-Wagstaff scenario of Proposition~\ref{prop:spike}. The one
property such models cannot imitate is exact divisibility: $q\mid2^p-1$ if and only if $2$ is a $\frac{q-1}p$-th power
residue modulo $q$, a residue condition of unbounded depth.

\subsection{Bounded-depth axioms for \texorpdfstring{$M_p$}{M\_p} and \texorpdfstring{$a_p$}{a\_p}}

\begin{definition}[Axioms $\mathrm{Ax}_{D,J}$]
A family assigns to each prime $p\ge5$ multisets of primes $F(p)$ (a fake $M_p$) and $G(p)$ (a fake $a_p$) such that:
\begin{enumerate}
\item[(A1)] $\prod F(p)$ has $p$ bits and $\prod G(p)$ has $p-1$ bits;
\item[(A2)] every element is $\equiv1\pmod{2p}$;
\item[(A3)] the mod-$8$ rules of Lemma~\ref{lem:mod8} and $q\equiv\pm1\pmod8$ for $F$-primes hold;
\item[(A4)] $\prod F\equiv M_p$ and $\prod G\equiv a_p\pmod{8p^J}$; this includes the Fermat-quotient congruence;
\item[(A5)] for $d\le D$ with $d\mid q-1$: $2^{(q-1)/d}\equiv1\pmod q$ if and only if $d\mid(q-1)/p$ for $F$-primes, or
$d\mid(q-1)/2p$ for $G$-primes;
\item[(A6)] fibres are disjoint, and $3\notin G(p)$;
\item[(A7)] repeated primes are Wieferich-type;
\item[(A8)] Brun--Titchmarsh holds, and (A9) Stewart's bound holds.
\end{enumerate}
\end{definition}

\begin{theorem}[Barrier; sketch, modulo Siegel--Walfisz]\label{thm:barrier}
The true factorizations satisfy every $\mathrm{Ax}_{D,J}$. For every $D,J$ there are families satisfying
$\mathrm{Ax}_{D,J}$ with $|G(p)|=1$ for all large $p$, that is, with ``$\tau(2^p+1)=4$'', and with either
(a) $|F(p)|=1$ for all large $p$, so $\beta_K=0$; or
(b) $|F(p)|\ge p^{1-\delta}/4-1$ for all large $p$, so $\beta_K=\infty$.
A diagonal family satisfies all $\mathrm{Ax}_{D,J}$ at once, eventually in $p$.
\end{theorem}
\begin{proof}[Construction]
$F(p)$ consists of $r$ primes ($r$ even) with $q_i\equiv1\pmod{2p}$, $q_i\equiv7\pmod8$, and
$q_i\equiv2\pmod\ell$ for every odd prime $\ell\le D$, chosen up to $\exp(p^\delta)$. A filler prime fixes (A1)
and (A4); because $r$ is even, the filler can be taken $\equiv7\pmod 8$. The condition $q\equiv2\pmod\ell$ means that
the only $d\le D$ dividing $q-1$ are $d\in\{1,2\}$, so (A5) reduces to the quadratic character of $2$, which (A3)
already fixes. $G(p)$ is a single prime in the right classes. The moduli are at most a fixed power of $\log x$, so
Siegel--Walfisz applies uniformly.
\end{proof}

\begin{proposition}[Sharpness]
(A4) with $8p^J>2^{p-1}$ forces $\prod F(p)=M_p$, and adding exact divisibility $\prod F\mid M_p$ recovers the truth.
So fake models die only at $p$-adic depth $J\gtrsim p\log2/\log p$, or at residue depth $(q-1)/p$, which is unbounded.
\end{proposition}

\begin{lemma}[Mod-24 limit of the construction]
No product of primes that are all $\equiv7\pmod 8$ and $\equiv2\pmod3$ is $\equiv M_p\pmod{24}$.
\end{lemma}
\begin{proof}
$M_p\equiv7\pmod8$ needs an odd number of factors, and $M_p\equiv1\pmod3$ needs an even number.
\end{proof}
So adding product congruences modulo $3$ forces genuine Chebotarev conditions in non-abelian Kummer fields such as
$\Q(\zeta_3,2^{1/3})$. The construction above does not handle these; that extension is \emph{unproved}.

\begin{theorem}[Barrier for the constant; sketch, using Bombieri--Vinogradov and Heath-Brown's Linnik bound~\cite{HB92}]
\label{thm:constbarrier}
Fix $\varepsilon\in(0,\frac1{10})$ and $D\ge3$. For all primes $p$ outside a set of relative density zero, there are a set
$Q'_p$ of at least $(\frac{\log2}2-\varepsilon)\,p/\log p$ primes in $[p^2,p^2(\log p)^{20}]$ and a prime $r$ with the
following properties. Each of them is $\equiv1\pmod{2p}$, $\equiv7\pmod8$, and $\not\equiv1\pmod\ell$ for odd $\ell\le D$.
The product $N'=r\prod Q'_p$ satisfies $N'<M_p$ and $N'\equiv M_p\pmod{2^ip^j}$ for all $i\le\varepsilon p/12$,
$j\le\varepsilon p/(12\log p)$. Brun--Titchmarsh and Stewart's bound hold for them. So no argument that uses only these
properties can improve the constant in Lemma~\ref{lem:trivial}.
\end{theorem}

Any improvement of $\frac{\log 2}{2}$ is equivalent to bounding the product of the $q\mid M_p$ in
$[p^{2-\eta},p^{2+\eta}]$ by $2^{(1-\delta)p}$. This is a concentration lemma proved in the working notes.
This window is Hooley's middle range in Artin's conjecture.

%======================================================================
\section{A genuine obstruction: honest ideals with many prime factors}
%======================================================================
The fake models violate exact divisibility. The next family does not: it consists of genuine elements of $\cO$.

\begin{theorem}\label{thm:honest}
Let $p\ge5$, $p\nmid k$, and $\alpha_k=2^k-\zeta$. Then:
\begin{enumerate}
\item $\Nm(\alpha_k)=\Phi_p(2^k)=\prod_{d\mid k}\Phi_{dp}(2)$ and $\cO/(\alpha_k)\cong\Z/\Phi_p(2^k)$;
\item for every prime $q\ne p$ dividing $\Nm(\alpha_k)$, $2^k$ has order exactly $p$ modulo $q$, and
$(\alpha_k)=\prod_q(q,\zeta-2^k)^{v_q(\Nm\alpha_k)}$, a product of degree-one primes, one above each $q$;
\item $\omega(\alpha_k)\ge\tau(k)$. Choosing $k<p$ with $\tau(k)$ maximal gives
$\log\omega(\alpha_k)\ge(\log2+o(1))\log p/\log\log p$.
\end{enumerate}
\end{theorem}
\begin{proof}
(1) $\Nm(2^k-\zeta)=\prod_{a}(2^k-\zeta^a)=\Phi_p(2^k)$. Since $p\nmid k$, $\Phi_p(x^k)=\prod_{d\mid k}\Phi_{dp}(x)$. Also
$\cO/(2^k-\zeta)\cong\Z[x]/(\Phi_p(x),x-2^k)\cong\Z/\Phi_p(2^k)$.
(2) If $q\mid\Phi_p(2^k)$ and $2^k\equiv1\pmod q$, then $\Phi_p(2^k)\equiv p\pmod q$, so $q=p$. Hence $2^k$ has order
$p$. Because the quotient is cyclic, exactly one prime above $q$ divides $(\alpha_k)$, namely $(q,\zeta-2^k)$, to the
exponent $v_q(\Nm\alpha_k)$.
(3) Each $\Phi_{dp}(2)$, $d\mid k$, has a primitive prime (P2), since $dp\notin\{1,6\}$. Distinct $d$ give distinct
primes. Finally, (P5) gives $\max_{k<p}\tau(k)=2^{(1+o(1))\log p/\log\log p}$.
\end{proof}

A direct computation for $p=13$, $k=12$ gives $\omega(\alpha_{12})=13\ge\tau(12)=6$.

\begin{remark}
The structural identities used by class-group, Stickelberger, Galois-orbit and successive-minima approaches hold for
$\alpha_k$ exactly as for $\alpha_1=2-\zeta$. Examples are the sub-ideal dictionary $\cO/I\cong\Z/d_I$, free Galois orbits of the primes, and
Stickelberger relations. The only datum separating $2-\zeta$ from $\alpha_k$ is quantitative: the house is $3$
against $2^k+1$, and $\log\Nm\approx p\log2$ against $kp\log2$. So a proof of Conjecture~\ref{conj:beta} must use the
\emph{size} of the generator in a way stronger than the norm bound. One caveat: $\alpha_k$ does not satisfy the analogue
of $p\in P_K$, since $\Nm(2^k+\zeta)$ has at least $\tau(\text{odd part of }k)$ prime factors. The coupling with
$2+\zeta$ enters the methods tested below only through the residue $2+\zeta\equiv4$, which needs no factorization.
\end{remark}

%======================================================================
\section{The function-field analogue is true}
%======================================================================
For $g\in\F_\ell[T]$, $\tau(g)$ counts monic divisors and $\omega(g)$ counts distinct monic irreducible factors.

\begin{theorem}\label{thm:ff}
Let $\ell$ be an odd prime, $a\in\F_\ell[T]$ of degree $d\ge1$, $p\ne\ell$ an odd prime, and $f=\ord_p\ell$. Then
\begin{enumerate}
\item $\omega(a^p-1)\le d\,(1+(p-1)/f)$;
\item $\omega(a^p+1)\ge(p-1)/f$;
\item if $\tau(a^p+1)\le K$, then $\omega(a^p-1)\le d\,(1+\log_2K)$, uniformly in $p$.
\end{enumerate}
\end{theorem}
\begin{proof}
Let $\mu_p\subset\F_{\ell^f}$ be the $p$-th roots of unity. The Frobenius $x\mapsto x^\ell$ acts on
$\mu_p\setminus\{1\}$ with $(p-1)/f$ orbits of size $f$. For $\eta\ne1$ let
$\Nm(a-\eta)=\prod_{\sigma}(a-\eta^\sigma)\in\F_\ell[T]$, the product over $\mathrm{Gal}(\F_{\ell^f}/\F_\ell)$.

(1) $a^p-1=(a-1)\prod_{\text{orbits}}\Nm(a-\eta)$. At most $d$ primes divide $a-1$. If an irreducible $\pi$ divides
$a^p-1$ but not $a-1$, then $a$ has order $p$ in $(\F_\ell[T]/\pi)^\times$, a cyclic group of order
$\ell^{\deg\pi}-1$. So $f\mid\deg\pi$. Each $\Nm(a-\eta)$ has degree $df$, so it has at most $d$ distinct
irreducible factors of degree $\ge f$.

(2) $a^p+1=\prod_{\eta\in\mu_p}(a+\eta)$, since $p$ is odd. The norms $\Nm(a+\eta)$ over distinct orbits are nonconstant.
They are pairwise coprime: a common prime factor in $\F_{\ell^f}[T]$ of $a+\eta^\sigma$ and $a+\eta'^{\sigma'}$ would
divide the nonzero constant $\eta^\sigma-\eta'^{\sigma'}$.

(3) $2^{(p-1)/f}\le2^{\omega(a^p+1)}\le\tau(a^p+1)\le K$, so $(p-1)/f\le\log_2K$; now apply (1).
\end{proof}

\begin{remark}[Why it does not transfer]
The constant-field Frobenius order $f$ controls both $a^p-1$ and $a^p+1$. Over $\Z$, every prime above $q\mid M_p$ and
above $Q\mid a_p$ has residue degree $1$ in $\Q(\zeta_p)$, so no shared invariant exists. Over $\Z$, a uniform bound
$\omega(M_p)\le C(K)$ on $P_K$ would be the literal analogue of (3). The data already strain it: $p=43,79,191,701$ have
$a_p$ prime and $\omega(M_p)=3,3,5,8$.
\end{remark}

%======================================================================
\section{Exact-divisibility tools, and the single open estimate}\label{sec:tools}
%======================================================================

\begin{lemma}[Incidence/norm lemma]\label{lem:incidence}
For $q\mid M_p$, the primes of $\cO$ above $q$ are $\mathfrak q_{q,j}=(q,\zeta-2^j)$, $1\le j\le p-1$. They are
distinct, of degree $1$, and $v_{\mathfrak q_{q,j}}(\zeta-2^j)=v_q(M_p)$. Consequently, for $P\in\Z[x]$ with
$\Phi_p\nmid P$,
\[
\sum_{j=1}^{p-1}\log\gcd\bigl(P(2^j),M_p\bigr)\ \le\ \log|\Nm(P(\zeta))|\ \le\ (p-1)\log\|P\|_1 .
\]
\end{lemma}
\begin{proof}
$2$ has order $p$ modulo $q$, so the roots of $\Phi_p$ modulo $q$ are $2^j$, $1\le j\le p-1$, and $q$ splits completely.
For $1\le j\le p-1$, $\Nm(2^j-\zeta)=\Phi_p(2^j)$. By (P3), $v_q(\Phi_p(2^j))=v_q(2^{jp}-1)=v_q(M_p)+v_q(j)=v_q(M_p)$.
The quotient $\cO/(2^j-\zeta)$ is cyclic, so only $\mathfrak q_{q,j}$ above $q$ divides $(2^j-\zeta)$, to exponent
$v_q(M_p)$. Hence $P(\zeta)\equiv P(2^j)\pmod{\mathfrak q_{q,j}^{\,v_q(M_p)}}$, and
$v_{\mathfrak q_{q,j}}(P(\zeta))\ge\min(v_q(P(2^j)),v_q(M_p))$. The ideals $\mathfrak q_{q,j}$ are distinct of norm
$q$, which gives the first inequality. The second is $|P(\zeta^a)|\le\|P\|_1$.
\end{proof}
In the fake models this fails: for a fake $q$, $(q,\zeta-2)=(1)$. It is the basic exact-divisibility input.
Lattice-existence (Siegel, Minkowski) arguments cannot extract a bound on $|S|$ from it, because both sides agree up to
factors independent of $|S|$.

\begin{lemma}[Divisor lattices]\label{lem:O}
Let $s>N^\theta$ with $\gcd(s,N)=1$. Let a nonempty set $D'$ of divisors of $N$ meet $R$ residue classes modulo $s$.
Then $\theta\,\omega(N)<1+\log_2R+\log_2(\tau(N)/|D'|)$.
\end{lemma}
\begin{proof}
Let $t=\lfloor\theta\,\omega(N)\rfloor$ and let $A$ be the product of the $t$ smallest prime-power parts of $N$. Then
$A\le N^{t/\omega(N)}\le N^\theta<s$. Write $N=AB$. Every divisor is $ab$ with $a\mid A$, $b\mid B$. If $ab\equiv a'b\pmod s$,
then $a\equiv a'$ because $b$ is invertible modulo $s$, so $a=a'$ because $a,a'\le A<s$. Hence each class contains at most
$\tau(B)=\tau(N)/\tau(A)\le\tau(N)/2^t$ divisors. So $|D'|\le R\,\tau(N)/2^t$, and $\theta\,\omega(N)<t+1$.
\end{proof}

A hypothesis that all divisors of $M_p$ lie in few classes modulo some $s\ge M_p^{1/4+\varepsilon}$ is what Lenstra's
theorem~\cite{Le84} or Coppersmith--Howgrave-Graham--Nagaraj~\cite{CHN08} would need. By Lemma~\ref{lem:O}, that
hypothesis already gives the conclusion up to a factor $4$ in $\omega$. So the input is the problem itself.

\begin{proposition}[$p$-adic constraints are void]
Write the prime factors of $M_p$, with multiplicity $w=\Omega(M_p)$, as $q_i=1+2pm_i$. Modulo $p^J$ the full constraint
is $\prod_i(1+2pm_i)\equiv M_p\pmod{p^{J+1}}$, one equation in the cyclic group $(1+p\Z)/(1+p^{J+1}\Z)$ of order
$p^J$. Its solution set in $(\Z/p^J)^w$ has exactly $p^{J(w-1)}$ elements, a density of $p^{-J}$ for every $w$. So no
number of factors is ever excluded, and the elementary symmetric functions of the $m_i$ are not individually determined.
\end{proposition}
\begin{proof}
$m\mapsto1+2pm$ is a bijection from $\Z/p^J$ onto that group, and $M_p\equiv1\pmod p$.
\end{proof}

\paragraph{Further tests.} The working notes test four more families of methods, each with a stated point of failure.
\begin{itemize}
\item \emph{Reciprocity.} Every power-reciprocity law between $2-\zeta$ and $2+\zeta$ reduces to local symbols needing
neither factorization. For $q\mid M_p$ and $n\mid q-1$ with $p\nmid n$, $2$ is automatically an $n$-th power.
\item \emph{Subgroup bounds.} $\langle2\rangle\subset\F_q^\times$ is \emph{the} subgroup of order $p$. So
Stepanov, Heath-Brown--Konyagin and Bourgain--Glibichuk--Konyagin bounds are identical for every $q\equiv1\pmod{2p}$.
\item \emph{Sparse multiples.} Pigeonhole forces $d\mid\sum_{i\le T}\pm2^{e_i}$ for every odd $d<\binom pT$. The
needed uniform bound on $\gcd(u,M_p)$ for sparse $u$ is false. For example, $\gcd(2+2^{30}-2^{11}-2^{12},M_{43})=9719$
with $43\in P_4$.
\item \emph{Kummer--Chebotarev.} $q\mid M_p$ if and only if $q$ splits completely in $\Q(\zeta_m,2^{1/m})$,
$m=(q-1)/p$. The degree grows like $(q/p)^2$, and each $m$ has one candidate $q$.
\end{itemize}

\paragraph{The single open estimate.} Every route above stops at the following statement:
\begin{equation}\label{eq:open}
\#\{q\le2^p:\ q\equiv1\ (2p),\ 2^p\equiv1\ (q)\}\le\exp\bigl(o(\log p/\log\log p)\bigr)\quad(p\in P_K\text{ or all }p).
\end{equation}
The best known bound is $(\frac{\log2}2+o(1))p/\log p$. We know of no pointwise bound that is even $o(p/\log p)$.
By Section~\ref{sec:barriers} and Theorem~\ref{thm:honest}, a proof of~\eqref{eq:open} must use exact divisibility of
unbounded depth \emph{and} the archimedean smallness of $2-\zeta_p$ beyond its norm.

%======================================================================
\section{Numerical data}\label{sec:data}
%======================================================================
All computations were redone independently for this paper. They use the factor-count tables
OEIS A046801 ($\tau(2^n-1)$, $n\le1206$), A085021 ($\omega(\Phi_n(2))$), A046800 ($\omega(2^n-1)$), A046798
($\tau(2^n+1)$, $n\le1128$) and A000978 (Wagstaff exponents)~\cite{OEIS}, together with direct computation. The
script \texttt{paper\_checks.py} reproduces every number below.

\begin{itemize}
\item $f(2n)\ge3f(n)$ holds for all $n\le606$, with equality only at $n=1,2,4$. Further,
$\min_{32\le n\le606}f(2n)/f(n)=14.256$, and $\min_{150\le n\le300}=256.1$.
\item $f(2n)/f(n)$ at $n=8,16,32,64,128$ is $3.741,\ 5.337,\ 15.134,\ 18.446,\ 154.843$. The minimum over
$16\le n\le150$ is at $n=16$. Local minima sit at powers of $2$, the configuration singled out by
Theorems~\ref{thm:three}--\ref{thm:fermat}.
\item The record values of $T(k)$ for $k\le600$ occur at
$k=1,2,4,6,8,12,18,20,24,36,48,60,72,84,108,120,144,168,180,288,300,360,420,540$.
\item $\max_{d<600}\omega(\Phi_d(2))=9$, attained only at $d=397$. Indeed
\[
M_{397}=2383\cdot6353\cdot50023\cdot53993\cdot202471\cdot5877983\cdot p_{30}\cdot p'_{31}\cdot p''_{31},
\]
where $p_{30}=814132872808522587940886856743$, $p'_{31}=1234904213576000272542841146073$ and
$p''_{31}=6597485910270326519900042655193$. We checked that the product is exactly $2^{397}-1$, that every factor is
$\equiv1\pmod{794}$, and that every factor passes the BPSW test. Also $\omega(M_{431})=8$.
\item The triples $(p,\tau(2^p+1),\omega(M_p))$ for $p\le293$ used in the working notes agree with the tables in all
61 cases. Every $\tau(2^p+1)$ with $5\le p\le1128$ is divisible by $4$ (Corollary~\ref{cor:4div}).
\item On $P_4$ with $p\le1128$, that is, on the Wagstaff exponents $5,\dots,701$, the values of $\omega(M_p)$ are
$1,1,2,1,1,1,2,1,3,1,3,2,1,2,5,2,4,2,8$. The maximum is $\omega(M_{701})=8$. On $P_8$ the maximum is also $8$.
\item \emph{Small factors versus Gillies.} For the 782 primes $3\le p\le6000$, let $c_p$ be the number of prime factors
$q\le p^2$ of $M_p$. The observed counts of $c_p=0,1,2,3,4$ are $436,264,65,16,1$; the maximum is $c_{397}=4$ and the
mean is $0.57$. Gillies' heuristic~\cite{Gi64} predicts $c_p$ is approximately Poisson with mean
$\mu_p=\log(\log p^2/\log(2p+1))$. The mixture of these Poisson laws predicts $428.8,\ 257.4,\ 77.5,\ 15.6,\ 2.4$.
An extremal configuration would need about $236$ such factors.
\end{itemize}

%======================================================================
\section{Corrections to the working notes and problem statements}\label{sec:errata}
%======================================================================
\begin{enumerate}
\item The round-2 note on the doubling problem states $\max_{d<600}\omega(\Phi_d(2))=8$ and
$\max_{p<600}\omega(2^p-1)=8$, at $p=431$. \textbf{Both are false}: $\omega(\Phi_{397}(2))=\omega(M_{397})=9$. The
problem statement for $\beta_K$ repeats the first claim for $d\le600$. The value $\omega(M_{431})=8$ is correct.
\item The note on $\beta_K$ says the heuristic predicts about $C\log\log x$ Wagstaff exponents up to $x$. The
prediction is $\asymp\log x$ exponents $p\le x$; the quantity $\log\log x$ is the count of Wagstaff \emph{numbers} up to
$x$. No proof is affected.
\item The proof of the no-certificate proposition in the $\beta_K$ note fails for $f=cx^k$; it is repaired in
Proposition~\ref{prop:nocert}.
\item The $\beta_K$ note's one-line summary of the constant barrier omits that the fake products satisfy $N'<M_p$
rather than having exactly $p$ bits. Theorem~\ref{thm:constbarrier} states it correctly.
\item The round-2 note says a negative answer to \#893 ``requires'' thin $k$ with $\tau(2^k+1)$ bounded. What is
proved unconditionally is Proposition~\ref{prop:negative}. Bounded $\tau(2^k+1)$ follows under the spike hypothesis of
Proposition~\ref{prop:spike}.
\item Problem-statement corrections recorded in the $\beta_K$ note, all confirmed. Brun--Titchmarsh is not needed for
Lemma~\ref{lem:trivial}. The $p$-adic digits of $M_p$ fix only $\sum_k(2p)^{k-1}e_k=(2^{p-1}-1)/p$, not each
elementary symmetric function $e_k$. ``$(2+\zeta_p)$ has at most $\log_2(K/2)$ prime-ideal factors'' holds for distinct
factors; with multiplicity the bound is $K/2-1$.
\end{enumerate}

%======================================================================
\section{Relation to the literature}\label{sec:novelty}
%======================================================================
A literature search was carried out in November 2026: erdosproblems.com, arXiv, the Kova\v{c}--Luca paper and its
references, and standard sources on Mersenne and Wagstaff numbers. ``New'' below means not found there.
\begin{center}
\small
\begin{tabular}{p{6.2cm}p{8.4cm}}
\toprule
Result & Status \\
\midrule
Identity (Thm.~\ref{thm:identity}) & Elementary; probably folklore. \\
$\limsup f(2n)/f(n)=\infty$ & \textbf{Known}: Kova\v{c}--Luca~\cite{KL}. Corollary~\ref{cor:KL} is a variant proof. \\
$f(2n)\ge3f(n)$, sharp (Thm.~\ref{thm:three}) & New; first $\liminf$ bound above $2$. Elementary. \\
Fermat bound, $\liminf\ge4$ if $F_a$ composite (Thm.~\ref{thm:fermat}) & New. Elementary. \\
Chains and density (Thm.~\ref{thm:density}) & New. Elementary. \\
Thin reduction and conditional result (Thm.~\ref{thm:thin}, Cor.~\ref{cor:cond}) & New. The hypothesis is strictly
weaker than the conjecture used in~\cite{KL}. \\
Necessary conditions; spike $\Rightarrow\beta\ge\log2$ (Props.~\ref{prop:negative}, \ref{prop:spike}) & New. \\
Soft-property and bounded-depth barriers (Sec.~\ref{sec:barriers}) & New as stated; proved in sketch. The method
(random or fake models) is standard in spirit, and the analytic input is classical. \\
Mod-8 structure, $4\mid\tau(2^p+1)$ & Elementary; the mod-8 facts on Wagstaff divisors are classical. \\
Trivial bound (Lemma~\ref{lem:trivial}) & Folklore. \\
Disproof $\Rightarrow$ open problems & Elementary implications of known open problems~\cite{Guy}. \\
Honest-ideal obstruction (Thm.~\ref{thm:honest}) & New framing; elementary. \\
Function-field analogue (Thm.~\ref{thm:ff}) & New as stated; an elementary consequence of the factorization of
cyclotomic polynomials over finite fields. \\
Divisor-lattice inequality (Lemma~\ref{lem:O}) & New to our knowledge; elementary; compare~\cite{Le84}. \\
Incidence lemma (Lemma~\ref{lem:incidence}) & Standard algebraic number theory. \\
Factor data & \textbf{Known}: Cunningham project~\cite{Cunningham}, OEIS. We contribute verification and corrections. \\
Conjecture~\ref{conj:beta} & Not found in the literature. Consistent with Gillies' heuristic. \\
\bottomrule
\end{tabular}
\end{center}

\emph{No result in this paper resolves Erd\H{o}s Problem~\#893 or Conjecture~\ref{conj:beta}.}
The new results are modest and elementary in technique. Their value lies in sharp constants, exact reductions and
precise statements of what any resolution must use.

\section{Open problems}
\begin{enumerate}
\item Is $\liminf f(2n)/f(n)>3$? Is it at least $4$ unconditionally, or $\infty$?
\item Prove~\eqref{eq:thinhyp}, or any bound $\omega(\Phi_d(2))\le\psi(d)^{1-\delta}$ on thin $d$.
\item Prove $\omega(M_p)=o(p/\log p)$ for all primes $p$, or even for a density-one set of $p$. By
Theorem~\ref{thm:constbarrier} this requires exact divisibility in Hooley's middle range.
\item Determine $\beta_K$. A first step would be to rule out $\beta_K\ge\log2$, which by Proposition~\ref{prop:spike}
would close the prime-spike route to a negative answer to \#893.
\end{enumerate}

\begin{thebibliography}{99}
\bibitem{Ba86} A.~S.~Bang, Taltheoretiske unders{\o}gelser, \emph{Tidsskrift Math.}~4 (1886), 70--80, 130--137.
\bibitem{BSW89} P.~T.~Bateman, J.~L.~Selfridge and S.~S.~Wagstaff, Jr., The new Mersenne conjecture,
\emph{Amer. Math. Monthly}~96 (1989), 125--128.
\bibitem{Cunningham} J.~Brillhart, D.~H.~Lehmer, J.~L.~Selfridge, B.~Tuckerman and S.~S.~Wagstaff, Jr.,
\emph{Factorizations of $b^n\pm1$}, 3rd ed., Contemp. Math.~22, AMS, 2002; and the Cunningham Project tables.
\bibitem{CHN08} D.~Coppersmith, N.~Howgrave-Graham and S.~V.~Nagaraj, Divisors in residue classes, constructively,
\emph{Math. Comp.}~77 (2008), 531--545.
\bibitem{EP} T.~F.~Bloom, Erd\H{o}s Problems, Problem \#893, \url{https://www.erdosproblems.com/893}.
\bibitem{Er98} P.~Erd\H{o}s, Some of my new and almost new problems and results in combinatorial number theory,
in \emph{Number Theory (Eger, 1996)}, de Gruyter, 1998, 169--180.
\bibitem{Gi64} D.~B.~Gillies, Three new Mersenne primes and a statistical theory, \emph{Math. Comp.}~18 (1964), 93--97.
\bibitem{Guy} R.~K.~Guy, \emph{Unsolved Problems in Number Theory}, 3rd ed., Springer, 2004.
\bibitem{HB92} D.~R.~Heath-Brown, Zero-free regions for Dirichlet $L$-functions, and the least prime in an arithmetic
progression, \emph{Proc. London Math. Soc.}~64 (1992), 265--338.
\bibitem{KL} V.~Kova\v{c} and F.~Luca, On the number of divisors of Mersenne numbers, arXiv:2506.04883 (2025);
\emph{Experimental Mathematics} (2026), doi:10.1080/10586458.2026.2636060.
\bibitem{Le84} H.~W.~Lenstra, Jr., Divisors in residue classes, \emph{Math. Comp.}~42 (1984), 331--340.
\bibitem{MV73} H.~L.~Montgomery and R.~C.~Vaughan, The large sieve, \emph{Mathematika}~20 (1973), 119--134.
\bibitem{OEIS} OEIS Foundation, The On-Line Encyclopedia of Integer Sequences, sequences A000978, A046798, A046800,
A046801, A085021, \url{https://oeis.org}.
\bibitem{St13} C.~L.~Stewart, On divisors of Lucas and Lehmer numbers, \emph{Acta Math.}~211 (2013), 291--314.
\bibitem{Wa83} S.~S.~Wagstaff, Jr., Divisors of Mersenne numbers, \emph{Math. Comp.}~40 (1983), 385--397.
\bibitem{Wi07} S.~Wigert, Sur l'ordre de grandeur du nombre des diviseurs d'un entier, \emph{Ark. Mat. Astron. Fys.}~3
(1907), 1--9.
\bibitem{Zs92} K.~Zsigmondy, Zur Theorie der Potenzreste, \emph{Monatsh. Math. Phys.}~3 (1892), 265--284.
\end{thebibliography}

\end{document}
